\section{Preliminaries}
\label{sec:preliminaries}

\subsection{Notation and Symbols}
The following is a brief summary of all symbols used including both english and german definitions.

\subsubsection*{Gear Geometry}
\begin{tabular}{@{}lll@{}}
    $m$         & Module                        & \textit{Modul} \\
    $m_n$       & Normal module                 & \textit{Normalmodul} \\
    $m_t$       & Transverse module             & \textit{Axialmodul} \\
    $z$         & Number of teeth               & \textit{Zähnezahl} \\
    $\alpha_t$  & Pressure angle (transverse)   & \textit{Eingriffswinkel} \\
    $x$         & Profile shift coefficient & \textit{Profilverschiebungsfaktor} \\
    $h_a^*$     & Addendum coefficient          & \textit{Kopfhöhenfaktor} \\
    $c^*$       & Tip clearance factor          & \textit{Kopfspielfaktor} \\
    $c$         & Tip clearance                 & \textit{Kopfspiel} \\
    $p$         & Pith                          & \textit{Teilung} \\
    $\alpha_n$  & Normal pressure angle         & \textit{Normaleingriffswinkel} \\
    $\alpha_t$  & Transverse pressure angle     & \textit{Stirneingriffswinkel} \\
    $\beta$     & Helix angle                   & \textit{Schrägungswinkel} \\
    $\beta_b$   & Base helix angle              & \textit{Grundschrägungswinkel} \\
    $\gamma$    & Lead angle at reference cylinder  & \textit{Steigungswinkel auf dem Teilzylinder} \\
\end{tabular}

\subsubsection*{Diameters}
\begin{tabular}{@{}lll@{}}
    $d_p$   & Pitch (reference) circle diameter ($d_p = m \cdot z$)
            & \textit{Teilkreisdurchmesser} \\
    $d_b$   & Base circle diameter ($d_b = d_p \cos\alpha_t$)
            & \textit{Grundkreisdurchmesser} \\
    $d_a$   & Addendum (tip) circle diameter
            & \textit{Kopfkreisdurchmesser} \\
    $d_f$   & Dedendum (root) circle diameter
            & \textit{Fusskreisdurchmesser} \\
    $d^*$   & Arbitrary diameter
            & -- \\
    $h_a$   & Addendum from reference pitch circle
            & \textit{Zahnkopfhöhe} \\
    $h_f$   & Dedendum from reference circle
            & \textit{Zahnfusshöhe}
\end{tabular}

\subsubsection*{Tooth Dimensions}
\begin{tabular}{@{}ll@{}}
    $s_0$   & Tooth thickness at pitch circle \\
    $s_b$   & Tooth thickness at base circle \\
    $\gamma$& Half tooth angle at base circle \\
\end{tabular}

\subsubsection*{Curve Parameters and Angles}
\begin{tabular}{@{}ll@{}}
    $\phi$      & Curve parameter (rolling angle) \\
    $\theta$    & Angle with respect to x-axis \\
    $\phi_0, \phi_{start}$    & Starting parameter value \\
    $\phi_{end}$    & Ending parameter value \\
\end{tabular}

\subsubsection*{Vectors and Functions}
\begin{tabular}{@{}ll@{}}
    $\vb{r}_{\text{inv}}(\phi)$         & Involute position vector \\
    $\vb{r}_{\text{undercut}}(\phi)$    & Undercut curve position vector \\
    $\vb{T}(\phi)$                      & Unit tangent vector \\
    $\vb{N}(\phi)$                      & Unit normal vector \\
    $\vb{R}(\theta)$                    & 2D rotation matrix \\
\end{tabular}

\subsection{Basic Gear Geometry}

The fundamental relationships between the gear parameters are summarized below. All formulas follow the (DIN) ISO 21771 standard.

\subsubsection*{Helical Parameter Conversions}
\label{subsec:helical_conversions}
The standardised inputs for an involute gear are the normal pressure angle $\alpha_n$ (typically $\SI{20}{\degree}$) and the normal module $m_n$, both defined in the section perpendicular to the tooth direction. For a helical gear with helix angle $\beta$, the construction of the tooth profile takes place in the transverse plane (perpendicular to the gear axis), which uses the transverse pressure angle $\alpha_t$ and transverse module $m_t$. The two pairs are related by:
\begin{align}
    \tan(\alpha_t) &= \frac{\tan(\alpha_n)}{\cos(\beta)} \\
    m_t &= \frac{m_n}{\cos(\beta)} \\
    \tan(\beta_b) &= \tan(\beta) \cdot \cos(\alpha_t)
\end{align}
The base helix angle $\beta_b$ describes the helix on the base cylinder, which is steeper than $\beta$ on the pitch cylinder. For spur gears ($\beta = 0$), the conversions collapse to $\alpha_t = \alpha_n$, $m_t = m_n$, and $\beta_b = 0$, so the transverse and normal sections coincide and the distinction can be ignored.

\subsubsection*{Circle Diameters}
Given the normal module $m_n$, number of teeth $z$, normal pressure angle $\alpha_n$, helix angle $\beta$, profile shift coefficient $x$, addendum coefficient $h_a^*$, and clearance coefficient $c^*$ (with $\alpha_t$ and $m_t$ derived as above):

\begin{align}
    d_p &= m_t \cdot z \\
    d_b &= d_p \cos(\alpha_t) \\
    d_a &= d_p + 2\,h_a = m_t \cdot z + 2\,m_n\cdot(h_a^* + x) \\
    d_f &= d_p - 2\,h_f = m_t \cdot z - 2\,m_n\cdot(h_a^* + c^* - x)
\end{align}

The addendum and dedendum heights depend on the profile shift and use the normal module since profile shift is applied perpendicular to the tooth direction:
\begin{align}
    h_a &= (h_a^* + x) \cdot m_n \\
    h_f &= (h_a^* + c^* - x) \cdot m_n
\end{align}

For spur gears ($\beta = 0$, so $m_t = m_n = m$), the diameter formulas collapse to the familiar single-module form $d_a = m\cdot(z + 2\,x + 2\,h_a^*)$ and $d_f = m\cdot(z + 2\,x - 2\,(h_a^*+c^*))$.

For standard involute gears:
\begin{center}
\begin{tabular}{@{}lll@{}}
    $\alpha_n = \SI{20}{\degree}$ & Normal pressure angle & \textit{Normaleingriffswinkel} \\
    $h_a^* = 1$ & Addendum coefficient & \textit{Kopfhöhenfaktor} \\
    $c^* = 0.25$ & Clearance coefficient & \textit{Kopfspielfaktor} \\
\end{tabular}
\end{center}

\subsection{Sign Convention}
The symbol $\mypm$ indicates a sign that depends on which flank is being described:
\begin{itemize}
    \item $\textcolor{red}{+}$: Right flank (counterclockwise rotating involute)
    \item $\textcolor{red}{-}$: Left flank (clockwise rotating involute)
\end{itemize}

The 2D rotation matrix $\vb{R}(\theta)$ and translation operations used throughout this document are defined in Section~\ref{sec:mathematical_tools}.

